\documentclass[11pt, a4paper]{article}
\usepackage[a4paper, margin=2.2cm]{geometry}
\usepackage{fontspec}
\usepackage[spanish, english, bidi=basic, provide=*]{babel}
\babelprovide[import, onchar=ids fonts]{spanish}
\babelprovide[import, onchar=ids fonts]{english}
\babelfont{rm}{Noto Sans}
\usepackage{amsmath, amssymb, amsthm}
\usepackage{enumitem}
\setlist[itemize]{label=-}
\usepackage{hyperref}

\theoremstyle{definition}
\newtheorem{definition}{Definición}[section]

\theoremstyle{plain}
\newtheorem{lemma}{Lema}[section]
\newtheorem{theorem}{Teorema}[section]
\newtheorem{corollary}{Corolario}[section]

\setlength{\parindent}{0pt}
\setlength{\parskip}{5pt}

\begin{document}

\begin{center}
    {\Large \textbf{UNIVERSIDAD NACIONAL DE COLOMBIA -- SEDE BOGOTÁ}} \\[4pt]
    \textbf{Matemáticas Discretas I (2025963) -- Documento Matemático Formal} \\[2pt]
    \textit{Proyecto: Paradoja Dimensional: Juego de Supervivencia y Optimización sobre Árboles Binarios de Decisión} \\[2pt]
    \small Diciembre de 2026
\end{center}

\hrule
\vspace{8pt}

\textbf{Integrantes del Grupo:}
\begin{itemize}[noitemsep, topsep=2pt]
    \item Sebastián Eisneider Rey Pretel
    \item \textit{[Nombre Completo Integrante 2]}
    \item \textit{[Nombre Completo Integrante 3]}
    \item \textit{[Nombre Completo Integrante 4]}
\end{itemize}

\section{Definiciones Formales del Modelo}

\begin{definition}[Árbol Binario de Decisión con Raíz]
El espacio de estados y decisiones del videojuego se define formalmente como una tupla $T = (V, E, r)$, donde:
\begin{enumerate}[noitemsep, topsep=2pt]
    \item $V$ es un conjunto finito y no vacío de vértices (escenarios narrativos), particionado en dos subconjuntos disjuntos:
    \[
        V = V_I \cup L, \qquad V_I \cap L = \emptyset
    \]
    donde $V_I$ denota el conjunto de \textbf{vértices internos} (estados de decisión) y $L$ denota el conjunto de \textbf{hojas} (estados terminales).
    \item $r \in V_I$ es un vértice distinguido denominado \textbf{raíz}, que representa el estado inicial de toda partida.
    \item $E \subseteq V_I \times V$ es el conjunto de aristas dirigidas (transiciones), generado por una función total de transición binaria:
    \[
        \delta: V_I \times \Sigma \to V, \qquad \Sigma = \{0, 1\}
    \]
    tal que para todo $u \in V_I$, las imágenes $u_0 = \delta(u, 0)$ y $u_1 = \delta(u, 1)$ satisfacen $u_0 \neq u_1$ (todo nodo interno tiene grado de salida exacto $\deg^+(u) = 2$).
    \item \textbf{Propiedad de árbol con raíz:} Todo vértice $v \in V \setminus \{r\}$ tiene grado de entrada exacto $\deg^-(v) = 1$, la raíz cumple $\deg^-(r) = 0$, y para cada $v \in V$ existe un único camino simple dirigido desde $r$ hasta $v$.
\end{enumerate}
\end{definition}

\begin{definition}[Profundidad, Altura y Subárboles Inducidos]
Para el árbol $T = (V, E, r)$ se definen las siguientes métricas estructurales:
\begin{enumerate}[noitemsep, topsep=2pt]
    \item \textbf{Profundidad $d(v)$:} Para cada $v \in V$, $d(v) \in \mathbb{N}$ es la longitud (número de aristas) del único camino dirigido desde $r$ hasta $v$. En particular, $d(r) = 0$ y si $w = \delta(u, a)$, entonces $d(w) = d(u) + 1$.
    \item \textbf{Altura de un vértice $h(v)$:} Se define recursivamente mediante:
    \[
        h(v) = \begin{cases}
            0 & \text{si } v \in L, \\
            1 + \max\big(h(\delta(v, 0)),\, h(\delta(v, 1))\big) & \text{si } v \in V_I.
        \end{cases}
    \]
    \item \textbf{Altura del árbol $h(T)$:} Es la altura de su vértice raíz, $h = h(T) = h(r) = \max_{v \in L} d(v)$.
    \item \textbf{Subárbol inducido $T_v$:} Para cada $v \in V$, $T_v = (V_v, E_v, v)$ es el subárbol binario con raíz en $v$ conformado por $v$ y todos sus descendientes alcanzables mediante iteraciones de $\delta$. Se denota por $L_v = L \cap V_v$ al conjunto de hojas de $T_v$.
\end{enumerate}
\end{definition}

\begin{definition}[Vector de Estado y Lógica de Supervivencia]
A cada escenario $v \in V$ se le asocia un \textbf{vector de estado discreto}:
\[
    S(v) = (e_v, p_v) \in \mathbb{N} \times \{0, 1\}
\]
donde $e_v \in \mathbb{N}$ representa las unidades de energía restantes en la pistola de portales y $p_v \in \{0, 1\}$ es la valuación de verdad de la proposición atómica $P(v)$: \textit{``el jugador posee fluido dimensional estable en el escenario $v$''}.

Definimos sobre $V$ la fórmula proposicional de \textbf{condición de supervivencia} $\Phi(v)$:
\[
    \Phi(v) \equiv (e_v > 0) \land (p_v = 1)
\]
Con base en ella, se establece la \textbf{función de evaluación terminal} $W: L \to \{0, 1\}$:
\[
    W(l) = \begin{cases}
        1 & \text{si } \Phi(l) \text{ es verdadera (Victoria)}, \\
        0 & \text{si } \neg\Phi(l) \text{ es verdadera (Derrota)}.
    \end{cases}
\]
\end{definition}

\begin{definition}[Partida y Conjunto de Estados Ganadores]
\begin{enumerate}[noitemsep, topsep=2pt]
    \item \textbf{Partida:} Una partida completa es una sucesión finita de vértices $\pi = (v_0, v_1, \dots, v_k)$ tal que $v_0 = r$, $v_k \in L$, y para cada $i \in \{0, \dots, k-1\}$ existe $a_i \in \{0, 1\}$ tal que $v_{i+1} = \delta(v_i, a_i)$. La palabra binaria asociada $\omega = a_0 a_1 \dots a_{k-1} \in \Sigma^k$ constituye la estrategia del jugador.
    \item \textbf{Conjunto de Estados Ganadores $G \subseteq V$:} Un vértice $v \in V$ pertenece a $G$ si y solo si existe al menos una hoja $l \in L_v$ en el subárbol $T_v$ tal que $W(l) = 1$.
\end{enumerate}
\end{definition}

\section{Lema Auxiliar sobre Subárboles Binarios}

\begin{lemma}[Descomposición y Altura de Subárboles Hijos]
Sea $T = (V, E, r)$ un árbol binario de decisión y sea $v \in V_I$ un vértice interno con hijos $v_0 = \delta(v, 0)$ y $v_1 = \delta(v, 1)$. Entonces:
\begin{enumerate}[noitemsep, topsep=2pt]
    \item Los conjuntos de vértices $V_{v_0}$ y $V_{v_1}$ son disjuntos ($V_{v_0} \cap V_{v_1} = \emptyset$), y el conjunto de hojas de $T_v$ es la unión disjunta $L_v = L_{v_0} \sqcup L_{v_1}$.
    \item La altura de ambos subárboles hijos satisface $h(v_0) \leq h(v) - 1$ y $h(v_1) \leq h(v) - 1$.
\end{enumerate}
\end{lemma}

\begin{proof}
Para (1), supóngase por contradicción que existe un vértice $z \in V_{v_0} \cap V_{v_1}$. Entonces existe un camino dirigido de $v_0$ a $z$ y otro de $v_1$ a $z$. Al concatenarlos con las aristas $(v, v_0)$ y $(v, v_1)$, se obtienen dos caminos dirigidos distintos desde $v$ hasta $z$ (ya que $v_0 \neq v_1$ por la Definición 1.1). Concatenando el único camino desde $r$ hasta $v$, existirían dos caminos distintos de $r$ a $z$, lo cual contradice la unicidad de caminos en un árbol con raíz. Por tanto, $V_{v_0} \cap V_{v_1} = \emptyset$, y al intersectar ambos conjuntos con $L$ resulta $L_v = L_{v_0} \sqcup L_{v_1}$.

Para (2), por la Definición 1.2 se tiene $h(v) = 1 + \max(h(v_0), h(v_1))$. Despejando, $\max(h(v_0), h(v_1)) = h(v) - 1$, de donde se sigue inmediatamente que $h(v_0) \leq h(v) - 1$ y $h(v_1) \leq h(v) - 1$.
\end{proof}

\section{Teoremas Principales y Demostraciones}

\begin{theorem}[Acotamiento de Estados Terminales y Terminación Garantizada]
Sea $T = (V, E, r)$ un árbol binario de decisión de altura $h \geq 0$. Entonces se cumplen las siguientes dos propiedades:
\begin{enumerate}[noitemsep, topsep=2pt]
    \item \textbf{Cota superior de hojas:} El número total de estados terminales está acotado por $|L| \leq 2^h$.
    \item \textbf{Terminación finita:} Toda partida $\pi = (v_0, v_1, \dots, v_k)$ finaliza en a lo sumo $k \leq h$ decisiones.
\end{enumerate}
\end{theorem}

\begin{proof}
Demostraremos la parte (1) mediante \textbf{inducción fuerte sobre la altura $h(v)$} de cualquier subárbol $T_v$, probando que $|L_v| \leq 2^{h(v)}$ para todo $v \in V$. El resultado global se obtendrá al evaluar en $v = r$.

\textbf{Caso base ($h(v) = 0$):} Si $h(v) = 0$, por la Definición 1.2 se tiene necesariamente que $v \in L$ (pues todo vértice interno cumple $h(v) \geq 1$). En consecuencia, $V_v = \{v\}$ y $L_v = \{v\}$, por lo que:
\[
    |L_v| = 1 = 2^0 = 2^{h(v)}.
\]
Así, la desigualdad se cumple con igualdad en el caso base.

\textbf{Hipótesis inductiva:} Sea $m \geq 1$ un entero. Supóngase que para todo vértice $u \in V$ cuya altura satisfaga $0 \leq h(u) < m$, se cumple que $|L_u| \leq 2^{h(u)}$.

\textbf{Paso inductivo ($h(v) = m$):} Sea $v \in V$ un vértice con $h(v) = m \geq 1$. Dado que $h(v) > 0$, se sigue que $v \in V_I$, por lo que $v$ posee exactamente dos hijos distintos $v_0 = \delta(v, 0)$ y $v_1 = \delta(v, 1)$. Por el Lema 2.1, se cumple que $L_v = L_{v_0} \sqcup L_{v_1}$ con $h(v_0) \leq m - 1 < m$ y $h(v_1) \leq m - 1 < m$. Aplicando el principio de adición para conjuntos disjuntos y la hipótesis inductiva sobre $v_0$ y $v_1$:
\[
    |L_v| = |L_{v_0}| + |L_{v_1}| \leq 2^{h(v_0)} + 2^{h(v_1)}.
\]
Dado que la función exponencial $f(x) = 2^x$ es estrictamente creciente en $\mathbb{N}$ y $h(v_0), h(v_1) \leq m - 1$, se obtiene:
\[
    |L_v| \leq 2^{m - 1} + 2^{m - 1} = 2 \cdot 2^{m - 1} = 2^m = 2^{h(v)}.
\]
Por el principio de inducción matemática, $|L_v| \leq 2^{h(v)}$ para todo $v \in V$. Evaluando en la raíz $v = r$ con $h(r) = h$, concluimos que $|L| \leq 2^h$.

Para la parte (2), sea $\pi = (v_0, v_1, \dots, v_k)$ una partida cualquiera con $v_0 = r$ y $v_k \in L$. Por construcción, $\pi$ es un camino simple dirigido de longitud $k$ desde $r$ hasta $v_k$, por lo que $k = d(v_k)$. Dado que $v_k \in L$ y $h = \max_{l \in L} d(l)$, se tiene:
\[
    k = d(v_k) \leq \max_{l \in L} d(l) = h.
\]
Por tanto, toda ejecución termina en a lo sumo $h$ decisiones.
\end{proof}

\begin{theorem}[Caracterización Recursiva del Conjunto de Estados Ganadores]
Sea $T = (V, E, r)$ un árbol binario de decisión y sea $G \subseteq V$ el conjunto de estados ganadores. Para todo vértice $v \in V$, se cumple la equivalencia lógica:
\[
    v \in G \iff \big(v \in L \land W(v) = 1\big) \lor \Big(v \in V_I \land \big(\delta(v, 0) \in G \lor \delta(v, 1) \in G\big)\Big).
\]
\end{theorem}

\begin{proof}
Sea $v \in V$ arbitrario. Dado que $\{V_I, L\}$ es una partición de $V$ (Definición 1.1), analizamos ambos casos disjuntos:

\textbf{Caso 1: $v \in L$ (Estado terminal).} Cuando $v \in L$, la proposición $v \in V_I$ es falsa, por lo que el lado derecho de la equivalencia se reduce lógicamente a $W(v) = 1$. Por otra parte, al ser $v$ una hoja, el único vértice en el subárbol $T_v$ es el propio $v$, luego $L_v = \{v\}$. Según la Definición 1.4:
\[
    v \in G \iff \exists l \in L_v \text{ tal que } W(l) = 1 \iff W(v) = 1.
\]
Así, la equivalencia se cumple cuando $v \in L$.

\textbf{Caso 2: $v \in V_I$ (Estado de decisión interna).} Cuando $v \in V_I$, la proposición $v \in L$ es falsa, reduciendo el lado derecho de la equivalencia a $\delta(v, 0) \in G \lor \delta(v, 1) \in G$. Sean $v_0 = \delta(v, 0)$ y $v_1 = \delta(v, 1)$. Por el Lema 2.1, $L_v = L_{v_0} \cup L_{v_1}$. Aplicando la Definición 1.4 y la distributividad del cuantificador existencial sobre la unión de conjuntos:
\begin{align*}
    v \in G &\iff \exists l \in L_v \, \big(W(l) = 1\big) \\
            &\iff \exists l \in (L_{v_0} \cup L_{v_1}) \, \big(W(l) = 1\big) \\
            &\iff \big(\exists l \in L_{v_0} \, (W(l) = 1)\big) \lor \big(\exists l \in L_{v_1} \, (W(l) = 1)\big) \\
            &\iff (v_0 \in G) \lor (v_1 \in G) \\
            &\iff \delta(v, 0) \in G \lor \delta(v, 1) \in G.
\end{align*}
Como ambos casos agotan el universo $V = V_I \cup L$, la equivalencia queda demostrada para todo $v \in V$.
\end{proof}

\begin{corollary}[Oráculo Exacto de Conteo de Rutas Ganadoras]
Para cada vértice $v \in V$, defínase la función $\eta: V \to \mathbb{N}$ como $\eta(v) = |\{l \in L_v \mid W(l) = 1\}|$. Entonces:
\[
    \eta(v) = \begin{cases}
        W(v) & \text{si } v \in L, \\
        \eta(\delta(v, 0)) + \eta(\delta(v, 1)) & \text{si } v \in V_I.
    \end{cases}
\]
Además, $v \in G \iff \eta(v) \geq 1$, y el algoritmo que evalúa $\eta(r)$ tiene complejidad temporal lineal $\mathcal{O}(|V|) = \mathcal{O}(2^h)$.
\end{corollary}

\begin{proof}
Si $v \in L$, entonces $L_v = \{v\}$ y el conjunto $\{l \in \{v\} \mid W(l) = 1\}$ tiene cardinalidad $1$ si $W(v) = 1$ y $0$ si $W(v) = 0$, es decir, $\eta(v) = W(v)$. Si $v \in V_I$, por el Lema 2.1 se tiene $L_{v_0} \cap L_{v_1} = \emptyset$, por lo cual $\{l \in L_v \mid W(l) = 1\} = \{l \in L_{v_0} \mid W(l) = 1\} \sqcup \{l \in L_{v_1} \mid W(l) = 1\}$. Por el principio de adición para conjuntos disjuntos, $\eta(v) = \eta(v_0) + \eta(v_1)$. Finalmente, cada vértice $v \in V$ es procesado una única vez en tiempo $\mathcal{O}(1)$ mediante recorrido en post-orden, tomando tiempo total $\mathcal{O}(|V|)$, donde $|V| \leq 2^{h+1} - 1 \in \mathcal{O}(2^h)$.
\end{proof}

\end{document}